\documentclass[a4paper, 12pt, french]{report}

% ========================
% IMPORTATION DES PACKAGES
% ========================

\usepackage[T1]{fontenc}
\usepackage[utf8]{inputenc}
\usepackage[margin=2cm]{geometry}
\usepackage{lmodern}
\usepackage{theorem}
\usepackage{amsfonts, amsmath, amssymb}
\usepackage{babel}

\pagestyle{headings}

\ifpdf
\usepackage[pdftex=true, hyperindex=true, colorlinks=false]{hyperref}
\else
\usepackage[pdftex=true, hyperindex=true, colorlinks=true]{hyperref}
\fi

\renewcommand{\familydefault}{\sfdefault}

\theoremstyle{break}
\theoremheaderfont{\scshape}
\theorembodyfont{\upshape}

\newtheorem{defin}{Définition}[section]
\newtheorem{prop}[defin]{Proposition}
\newtheorem{theo}[defin]{Théorème}
\newtheorem{coro}[defin]{Corollaire}

\title{DIFFERENTIABILITÉ DANS UN BANACH - APPLICATION}
\author{}
\date{\today}

\begin{document}
\maketitle

\tableofcontents

\chapter{ESPACE NORME ET ESPACE DE BANACH}

\section{Espace vectoriel norme et espace metrique}
\begin{defin}[Espace vectoriel norme]
    Un espace vecoriel norme $E$ est un espace vectoriel dans $\mathbb{K}$ ($\mathbb{R}$ ou $\mathbb{C}$) muni de l'application $\| \cdot \| : E \rightarrow \mathbb{R}_{+}$ verifiant les axiomes suivant:
    \begin{enumerate}
        \item Positivite : $\| \cdot \| = 0$ ssi $x = 0$ pour $x \in E$,
        \item Homogeneite : $\| \lambda x \| = |\lambda| \|x\|$, pour $\lambda \in \mathbb{R}$ et $x \in E$,
        \item Inegalite triangulaire : $\| x + y \| \leq \|x\| + \|y\|$, pour $x, y \in E$
    \end{enumerate}
\end{defin}

\begin{defin}
    Un espace metrique $E$ est un espace non vide muni de la distance $d$ verifiant les axiomes suivant:
    \begin{enumerate}
        \item Separation : $d(x,y) = 0$ ssi $x=y$, pour $x, y \in E$,
        \item Symetrie : $d(x,y)=d(y,x)$,
        \item Inegalite triangulaire : $d(x,z) \leq d(x, y) + d(y,z)$.
    \end{enumerate}
\end{defin}

\subsection*{Quelques exemple de norme}
\begin{itemize}
    \item \textbf{Sur $\mathbb{R}^{n}$}:
    \begin{itemize}
        \item[.] $\|x\|_{1} := \displaystyle \sum_{i=0}^{n} |x_{i}|$
        \item[.] $\|x\|_{2} := \displaystyle (\sum_{i=0}^{n} |x_{i}|^{2})^{\frac{1}{2}}$
        \item[.] $\|x\|_{\infty} := \displaystyle max_{i \in \{1,2, \cdots, n\}} \{ |x_{i}|\}$
    \end{itemize}
    \item \textbf{Sur $l^{p}$} pour $p \in [1, +\infty[$\\
    On definit $l^{p} := \{x = (x_{n})_{n \in \mathbb{N}}\;  | \; \displaystyle \sum_{i=0}^{n} |x_{i}| < +\infty\}$ 
    \begin{itemize}
        \item[.] Si $p = \infty$, $\|x\| := sup_{i \in \mathbb{N}} |x_{i}|$
    \end{itemize}
    \item \textbf{Sur $C_{b}(X)$} \\
    $C_{b}(X) : \{f : X \rightarrow \mathbb{R} | \mbox{ f est continue et borne }\}$
    \begin{itemize}
        \item[.] $\|f\|_{\infty} := \displaystyle sup_{x \in X} {|f(x)|}$
    \end{itemize}
    \item \textbf{Sur C([a,b])} \\
    $C([a,b]) := \{ f : [a,b] \rightarrow \mathbb{R} | \mbox{ f est continue }\}$
    \begin{itemize}
        \item[.] $\|f\|_{1} := \displaystyle \int_{0}^{1} |f(x)|dx$
        \item[.] $\|f\|_{2} := \displaystyle (\int_{0}^{1} |f(x)|^{2}dx)^{\frac{1}{2}}$
        \item[.] $\|f\|_{\infty} := \displaystyle sup_{x \in [a,b]}(|f(x)|)$
    \end{itemize}
\end{itemize}

\section{Topologie}
\begin{defin}
    On dit que $O$ est un ouvert dans $X$, s'il existe $x \in X$ on peut trouver $r > 0$ telle que $B(x,r) \subset O$.
    On dit que $F$ est un ferme dans $x$ si son complememtaire est un ouvert c-a-d $X \backslash F$ est  un ouvert.
\end{defin}

\begin{prop}
    \begin{itemize}
        \item $\emptyset$ et $X$ sont a la fois des ouverts et ferme,
        \item La reunion des ouverts quelconque sont des ouverts,
        \item L'intersection des ouverts finie sont des ouverts,
        \item La reunion de ferme finie est un ferme,
        \item L'intersection de ferme quelconque est un ferme.
    \end{itemize}
\end{prop}

\begin{prop}
    \begin{itemize}
        \item $\dot{A} := \{ B(x,r) \subset A\}$,
        \item $\overline{A} := \{ B(x,r) \cap A \not= \emptyset\}$,
        \item $Fr(A) := \overline{A} \backslash \dot{A}$
    \end{itemize}
\end{prop}

\begin{prop}
    \begin{itemize}
        \item A est ouvert ssi $A=\dot{A}$,
        \item A est un ferme ssi $A=\overline{A}$,
        \item $\bar{A} = Fr(A) \cup \dot{A}$.
    \end{itemize}
\end{prop}

\section{Suite de Cauchy et completude}
\begin{prop}
    $(x_{n})_{n \in \mathbb{N}}$ est une suite de Cauchy ssi $ \displaystyle \lim_{n, m \rightarrow \infty} \| x_{n} - x_{m}\| = 0$.
\end{prop}

\begin{prop}
    Toute suite de cauchy est borne.
\end{prop}

\begin{prop}
    Si $(x_{n})_{n \in \mathbb{N}}$ est convergente, alors elle est une suite de Cauchy.\\ En general, la reciproque est fausse pour un espace vectoriel norme quelconque.
\end{prop}

\begin{prop}
    Toute espace vectoriel de dimension finie est un espace de Banach. En particulier $(\mathbb{R}, \| \cdot \|_{2})$ est un espae de Banch.
\end{prop}

\begin{prop}
    $(l^{p}, \| \cdot \|_{\infty})$ est un espace de Banach pour $1 \leq p< +\infty$.
\end{prop}

\section{Proprieté fondamentale}
\begin{prop}
    $\| \cdot \|_{1}, \| \cdot \|_{2}$ sont equivalente s'il existe deux constante $c_{1}$ et $c_{2}$ telle que: $$c_{1}\| \cdot \|_{1} \leq \| \cdot \|_{2} \leq c_{2}\| \cdot \|_{1}.$$
\end{prop}

\begin{theo}
    Toute les normes sont equivalent sur un espace vectorielle.
\end{theo}

\begin{prop}
    Soit $E, F$ deux espaces de Banach et $T: E \rightarrow F$ une application. Lasse 
    \begin{itemize}
        \item $T$ est continue sur $E$,
        \item $T$ est continue sur 0,
        \item $T$ est borne (c-a-d $\|T(x)\|_{F} \leq M \|x\|_{E}$)
    \end{itemize}
\end{prop}

\begin{prop}
    $\|f\|_{\infty} :=  sup_{\|x\| \leq 1} \frac{\|f(x)\|_{F}}{\|x\|_{E}}$
\end{prop}

\end{document}